\section{Populations and measures}\label{sec:populations}

Which flights, which stretches of flight, and how many of them?

\paragraph{Groups.} The labels are those of the note \emph{Group flights in free flight}
(\texttt{methods/group-identification}). Two pilots are in contact when one is on the path of the
other within $\delta=2$~min, $\varepsilon_h=500$~m horizontally and $\varepsilon_z=300$~m vertically
(headings compared only in glide) for at least $\tau^*=\tauStar$~min. The configurations of an instant
are the connected components of the contact graph, and a group keeps the same members over a maximal
interval. A contact means shared air within sight; it does not prove interaction.

\paragraph{States.} After its \emph{departure}, the first fix 3~km from the takeoff, every instant of a
flight is in a group or solo; before it the pilot counts as neither. Flights that never depart are
left out. We use paraglider flights logged every 10~s or finer, so that the labels exist at every
instant; the phases (climb, search, glide) exist on the flights logged every second.

\paragraph{Populations.} Three, all of paragliders.
\begin{itemize}
  \item $\Sl$, \emph{solo flights}: flights that never entered a group, the whole of their time after
        the departure. They may have had contacts shorter than $\tau^*$.
  \item $\G$, \emph{group time}: the stretches in a group of the flights that entered one at least
        once. A stretch stays in one group; $\G$ counts a flight once however many groups it joins.
  \item $\A$, \emph{the same flights alone}: the solo stretches of those same flights, a control with
        the same pilots, gliders, loggers and days as $\G$.
\end{itemize}
The control of Section~\ref{sec:route} is stricter than $\Sl$ and is defined there.

\paragraph{Measure.} A \emph{piece} is a stretch within one segment of the track, of at least
1250~s, on a 10-s grid of its clock. For a flight $f$ and a population, the time-averaged mean
squared displacement pools the origins of all its pieces in that population,
\begin{equation}\label{eq:tamsd}
  \overline{\delta^2_f}(\tau)=\frac{1}{N_f(\tau)}\sum_{t}\bigl\|\mathbf r(t+\tau)-\mathbf r(t)\bigr\|^2 ,
\end{equation}
with $t$ and $t+\tau$ in the same piece. The ensemble $\MSD(\tau)$ is its mean over the flights, one
weight per flight, and $H$ is half the slope of $\log\MSD$ against $\log\tau$ over the lags from 10 to
1000~s: $H=1$ for a straight line, $H=\tfrac12$ for a random walk. $\Hloc(\tau)$ is the same slope over the
lags within $\pm0.25$ decades of $\tau$. It is read at $\tau=100$~s, within one climb or glide, and at
$\tau=1000$~s, over one to two climb--glide cycles. Climbs and glides are summarised by the share of
time in climb, the share of searches that end in a climb, the climb rate, the glide speed and the
glide ratio, each a ratio of sums over the phase runs of a population. Bands are 5--95\% of bootstrap
draws that resample the pilots and the dependency blocks of the grouping.

\begin{table}
\centering\small
\begin{tabular}{@{}l r r r@{}}
\toprule
 & solo flights $\Sl$ & group time $\G$ & same flights alone $\A$ \\
\midrule
flights, $H$ and $\Hloc$ & 119\,129 & 33\,435 & 29\,611 \\
pieces of at least 1250~s & 125\,997 & 45\,476 & 47\,577 \\
flights, time in climb and searches & 80\,278 & 24\,099 & 23\,469 \\
flights, climb rate and glides & 115\,629 & 33\,406 & 30\,703 \\
\bottomrule
\end{tabular}

\caption{Members of the populations, paragliders. The flights of $\G$ and $\A$ are the same flights,
and $\A$ keeps those with a solo stretch of at least 1250~s.}
\label{tab:populations}
\end{table}

\paragraph{Sizes.} Table~\ref{tab:populations} counts the members. $\G$ has \PopFlightsG{} flights in
\PopPiecesG{} pieces and $\Sl$ has \PopFlightsS{} flights in \PopPiecesS{}. $\A$ has \PopFlightsA{} flights
in \PopPiecesA{} pieces, fewer flights than $\G$: the others have no solo stretch long enough. The phase
statistics rest on fewer flights, those logged every second.
